\paragraph{C.3l (Proved): the two universal abelian injectivity assertions have one Beauville gate.}
Fix an algebraically closed field $k$ of arbitrary characteristic.
All cycle groups in this paragraph have rational coefficients.
The following two assertions are equivalent:
\begin{enumerate}
\item For every abelian variety $A/k$, both injectivity assertions
      of C.3h hold in every indicated degree.
\item For every abelian variety $V/k$, the ordinary Chow group
      has no negative Beauville pieces:
      $\mathrm{CH}^r(V)_{(t)}=0$ for every $r$ and $t<0$.
\end{enumerate}
The second assertion is the usual Beauville vanishing conjecture,
recalled in \citet[Equation (5.9), p.~262]{rootBeauville2010}.
We prove the equivalence, not the conjecture. Its quantifier over
all abelian varieties is essential: a family on $A$ introduces
the Jacobian of its parameter curve.

\begin{proof}
The additional point beyond C.3h is the following parameterization.
For each $z\in J^r(A)_{(s)}$, there are finitely many smooth
projective pointed curves $C_j$, points $a_j\in J(C_j)(k)$,
and cycles
\[
 Z_j\in\mathrm{CH}^r(J(C_j)\times A)_{(s-1)}
 \quad\text{such that}\quad
 z=\sum_j\bigl((Z_j)_{a_j}-(Z_j)_0\bigr).
\]
Here the subscript on $Z_j$ denotes the Beauville grading for
simultaneous multiplication on both abelian factors; its
evaluation is the refined pullback to the indicated point.

To prove this, express ordinary algebraic equivalence by curve
families: $z$ is a finite sum of $W_{j*}\delta_j$, where
$W_j\in\mathrm{CH}^r(C_j\times A)$ and
$\delta_j\in\mathrm{CH}_0(C_j)^0_{\mathbb Q}$.
Clear the finitely many rational coefficients, incorporating
them in $W_j$, so that each $\delta_j$ is an integral
degree-zero divisor class. It determines a point
$a_j\in J(C_j)(k)=\operatorname{Pic}^0(C_j)(k)$.
Let $D_j=c_1(\mathcal P_j)$ for the Poincar\'e line bundle on
$J(C_j)\times C_j$, normalized along zero and the base point
of $C_j$. On $J(C_j)\times C_j\times A$ put
\[
 \widetilde Z_j=(p_{13})_*
               \bigl(p_{12}^*D_j\cdot p_{23}^*W_j\bigr).
\]
The codimension is $1+r-1=r$. The Poincar\'e bundle's
defining universal property and its rigidification give
$([m]\times1)^*D_j=mD_j$: tensor power $m$ classifies
multiplication by $m$ on $\operatorname{Pic}^0(C_j)$,
and the rigidification removes a pullback from the base.
Fix $m>1$ prime to the characteristic. Flat base change
therefore gives $([m]\times1)^*\widetilde Z_j=m\widetilde Z_j$.
Decompose $\widetilde Z_j$ by the ordinary Beauville grading
of the abelian variety $J(C_j)\times A$. Its projectors are
polynomials in simultaneous multiplication and commute with
the first-factor operator. Thus each piece $Z_{j,(t)}$
retains first-factor eigenvalue $m$, while its total
eigenvalue is $m^{2r-t}$. The first-factor pullback is
injective on rational Chow groups: its finite-flat
pushforward composed with pullback is its nonzero degree.
Cancelling this pullback consequently gives
\[
 (1\times[m])^*Z_{j,(t)}=m^{2r-t-1}Z_{j,(t)}.
\]
Its fibre at any point of $J(C_j)$ therefore has $A$-index
$t+1$. The unprojected fibre at $a_j$ is $W_{j*}\delta_j$,
and the zero fibre is zero, by normalization of the
Poincar\'e bundle. These statements follow from commutation
of proper pushforward with refined pullback in the product
square. Every graded piece also has zero fibre at zero:
evaluation there commutes with first-factor multiplication,
so this fibre equals $m$ times itself. Taking the $s$-piece
of the sum of the evaluated classes now proves the required
formula with $Z_j=Z_{j,(s-1)}$. This uses only the ordinary
grading on an abelian product, without a relative Chow
projector on $C_j\times A$. No characteristic-zero argument
or isogeny denominator is involved; rational coefficients
are used for the Beauville decomposition itself.

Assume (2). If $s\le0$, all of these parameterizing cycles
vanish, since $s-1<0$ on the abelian varieties
$J(C_j)\times A$. Hence $J^r(A)_{(s)}=0$ for $s\le0$.
The quotient $B^r(A)$ also has no negative piece, directly
from (2). C.3h now proves both assertions in (1).

Conversely (1) and C.3h give $B^r(A)_{(s)}=0$ for $s<0$
and $J^r(A)_{(s)}=0$ for $s\le0$. The degreewise exact
sequence for the invariant subgroup and its quotient is
\[
 0\longrightarrow J^r(A)_{(s)}
 \longrightarrow\mathrm{CH}^r(A)_{(s)}
 \longrightarrow B^r(A)_{(s)}\longrightarrow0.
\]
For $s<0$ both ends vanish, proving (2).
\end{proof}

This uniform equivalence does not replace the smaller,
fixed-dimension tests in C.3h by a proved vanishing theorem.
In particular, it supplies no dimension bound on the
Jacobians of parameter curves for a proposed class in
$J^2(A)_{(0)}$.
